\chapter{Типизированное Callias--\texorpdfstring{$M_4$}{M4} вложение
фермионного cross-оператора}
\label{ch:v10-callias-m4-cross-embedding}

Каллиасов носитель
\begin{equation}
 \mathcal H_C=\mathbb C^2_{\rm spin}\otimes\mathbb C^2_{\rm twist}\otimes H_{15}
\end{equation}
уже имеет размерность $60$. Пространственные Pauli-генераторы действуют на
первом факторе, а равнозарядные twist-генераторы --- на втором; две
тройки коммутируют.

Положив
\begin{equation}
 R=I_2\otimes\tau_z\otimes I_{15},\qquad
 X=I_2\otimes\tau_x\otimes I_{15},
\end{equation}
получаем $R^2=X^2=I_{60}$ и $RX+XR=0$. Любой заряд вида
$I_4\otimes Q_{15}$ коммутирует с $X$. Проекторы $(I\pm R)/2$ имеют ранги
$30$, а частичная изометрия $P_+XP_-$ --- ранг $30$.

Pauli tensor basis на spin--twist-факторе имеет Gram-матрицу $60I_{16}$,
поэтому вложение полной алгебры $M_4$ имеет ранг $16$. При единичном
cross-поле Dirac-оператор имеет спектр
$(-\sqrt2)^{30},(\sqrt2)^{30}$, а determinant-кривизна равна $-60$:
пятнадцатиканальная кратность точно усиливает прежний четырёхмерный ответ.

Условное усиление клеточной twist-двойки задаётся картой
\begin{equation}
 E=I_2\otimes(1,\ldots,1)^T,\qquad E^*E=15I_2,\qquad \rank E=2.
\end{equation}
Но унаследованная карта равна нулю и имеет ранг $0$. Следовательно,
математическое вложение существует, однако выбор одинакового действия на
всех пятнадцати каналах и отождествление клеточной двойки с twist-фактором
не следуют из текущего родителя.

\begin{theorem}[Условное типизированное вложение]
Callias carrier размерности $60$ поддерживает равнозарядный cross-оператор,
полную алгебру $M_4$ и пятнадцатикратно усиленную фермионную
восприимчивость без конфликта с пространственным Clifford action.
\end{theorem}

\begin{nogocriterion}[Равномерность по $H_{15}$ требует родителя]
Запись $A\mapsto A\otimes I_{15}$ является условной архитектурой, а не
унаследованным выводом, пока физический оператор не выбирает один и тот же
cross-coupling во всех пятнадцати каналах.
\end{nogocriterion}

\begin{protocol}\begin{enumerate}
\item размерность носителя: \textbf{$60$};
\item ранг вложенной $M_4$-алгебры: \textbf{$16$};
\item ранг cross-изометрии: \textbf{$30$};
\item charge defect: \textbf{rank $0$};
\item determinant-кривизна: \textbf{$-60$};
\item условная архитектура: \textbf{$16/16$};
\item унаследованные данные: \textbf{$3/5$};
\item физическое происхождение: \textbf{$0/3$};
\item следующий гейт: \textbf{родитель равномерного усиления по $H_{15}$}.
\end{enumerate}\end{protocol}

Машинный сертификат сохранён в
\path{s2t_v10_particle_wrinkle_dislocation_callias_equal_charge_twist_m4_fermionic_cross_typed_embedding_gate_results.json}.